% ============================================================
% Appendix --- Problem sheet index
%
% This appendix exists because the restructure dissolved the official problem sheets: each
% problem now sits with the material it tests, so a sheet's problems can be spread over
% several chapters. The tables below put the sheet view back, and they are also where
% Corsin's colour-coded "Recommended exercises" pages are preserved --- those were
% per-sheet artefacts with nowhere else to live once the \exercisesheet{N} headings went.
%
% Built up during the migration, one sheet per stage.
% ============================================================
\chapter{Problem Sheet Index}
\label{ch:sheet_index}

% Transition: Claude Opus 5 (Medium)
The problem sheets are not reproduced as blocks anywhere in this document. Every problem is
quoted verbatim from the official sheet and placed next to the material it tests, with its
solution in the \emph{Solutions} section of the chapter that hosts it. That is the right
arrangement for reading the notes and the wrong one for answering \qt{what is on sheet 7?}, so
this appendix restores the sheet view.

All statements are quoted from \texttt{exercises/ExN\_Analysis2\_eng.pdf}; attribution:
Prof.\ Joaquim Serra, D-MATH, ETH Z\"urich.

% Tutor references moved out of the typeset text 2026-11-04 (the book names tutors only in source
% citations): Corsin opens each week's file with the colour-coded priority page; the "Note" column
% below quotes his comments on it. The one printed citation is in the notation block.
\begin{notation}[Priority key]
\label{not:priority_key}
From sheet 4 onwards each problem carries a priority, taken from the colour-coded
\textit{Recommended exercises} pages of Corsin Nick's class notes: % tutor-citation
blue \textbf{= important},
orange \textbf{= semi-important}, red \textbf{= optional}. The official sheets additionally mark
harder problems with \textbf{(*)}. Each problem carries its tag in its own title throughout this
document; the tables below collect them, together with the notes quoted from the same pages.
\end{notation}

\begin{ainote}
Sheets 1--3 have no priority page, so no problem on them is singled out, and all of them are
reproduced and solved.
\end{ainote}

\section{Sheets 1--3}
\label{sec:sheet_index_1_3}

\begin{center}
\begin{tabular}{@{}l l l@{}}
\textbf{Problem} & \textbf{Topic} & \textbf{Now in} \\
\hline
1.1 & Recap questions          & \cref{sec:analysis1_recap} \\
1.2 & From equations to drawings & \cref{sec:visualizing_sets} \\
1.3 & From drawings to equations & \cref{sec:visualizing_sets} \\
1.4 & Young's inequality       & \cref{sec:youngs_inequality} \\[4pt]
2.1 & Metric spaces            & \cref{sec:metric_spaces} \\
2.2 & Multiple choice (closure) & \cref{sec:unit_balls} \\
2.3 & Multiple choice (distance from a set) & \cref{sec:unit_balls} \\
2.4 & Boundary, interior       & \cref{sec:unit_balls} \\
2.5 & Product of metric spaces & \cref{sec:metric_spaces} \\
2.6 & Continuity of the distance & \cref{sec:continuity} \\
2.7 & Continuity of the composition & \cref{sec:continuity} \\[4pt]
3.1 & Lipschitz or not         & \cref{sec:lipschitz_continuity} \\
3.2 & Problems at the origin   & \cref{sec:continuity} \\
3.3 & Open, closed, complete and compact & \cref{sec:why_compactness_matters} \\
3.4 & Multiple choice (open, complete, compact) & \cref{sec:why_compactness_matters} \\
3.5 & Multiple choice (fixed points) & \cref{sec:banach_fixed_point} \\
3.6 & Multiple choice (uniform continuity) & \cref{sec:lipschitz_continuity} \\
3.7 & Cauchy--Schwarz          & \cref{sec:cauchy_schwarz} \\
3.8 & Inner product and norm   & \cref{sec:cauchy_schwarz} \\
3.9 & All norms are equivalent & \cref{sec:cauchy_schwarz} \\
3.10 & Hilbert--Schmidt norm   & \cref{sec:cauchy_schwarz} \\
3.11 & (*) Uniform continuity and moduli & \cref{sec:lipschitz_continuity} \\
\end{tabular}
\end{center}

\section{Sheet 4}
\label{sec:sheet_index_4}

% Source: Corsin Nick/Class Notes/Week 4.pdf, p. 1
% Corsin's "Recommended exercises" page, preserved here: it was a per-sheet artefact and had
% nowhere else to go once the \exercisesheet{4} heading was dissolved.
\begin{center}
\begin{tabular}{@{}l p{0.28\linewidth} p{0.28\linewidth} l@{}}
\textbf{Problem} & \textbf{Priority} & \textbf{Note} & \textbf{Now in} \\
\hline
4.1 & important & ``Use the chain rule'' & \cref{sec:jacobi_matrix} \\[3pt]
4.2 & important & ``Recall that path connected implies connected'' & \cref{sec:path_connectedness} \\[3pt]
4.3 & semi-important (parts 1, 2, 3, 5); \textbf{parts 4 and 6 important} &
  ``For 3., use Cauchy--Schwarz in $\mathbb{R}^n$ applied to convenient vectors.'' &
  \cref{sec:cauchy_schwarz} \\[3pt]
4.4 & optional & --- & \cref{sec:cauchy_schwarz} \\[3pt]
4.5 & optional & --- & \cref{sec:mean_value_theorem} \\[3pt]
4.6 & important (parts 1, 2, 3) & --- & \cref{sec:chain_rule_along_path} \\
\end{tabular}
\end{center}

% Correction: Claude Opus 5.5 (Extra) --- the printed notes under sheets 4, 5, 6, 8, 12 and 13
% recorded the repository's own history ("appear for the first time", "added on 2026-08-11",
% "until this restructure", "considered and rejected"), which tells a reader of the PDF nothing
% and is what git log is for. Kept below as comments; the printed notes keep only what a reader
% can use.
% Problems 4.4 and 4.5 had been transcribed into content/exercise-sheets/week-04.tex, but that
% file was never \input by anything, so neither problem had ever reached the built PDF and
% neither carried a solution. Both are now placed and solved.

\section{Sheet 5}
\label{sec:sheet_index_5}

% Source: Corsin Nick/Class Notes/Week 5.pdf, p. 1
\begin{center}
\begin{tabular}{@{}l l p{0.42\linewidth} l@{}}
\textbf{Problem} & \textbf{Priority} & \textbf{Note} & \textbf{Now in} \\
\hline
5.1 & important & ``Recall the linearity of the differential and that $C^1 \Rightarrow$
  differentiable'' & \cref{sec:directional_derivatives} \\[3pt]
5.2 & important & ``Compare with the example from the script and recall that continuous
  functions on compact sets attain their extrema.'' & \cref{sec:lagrange_multipliers} \\[3pt]
5.3 & important & ``There is no difference between $x \to 0$ and $|x| \to 0$.'' &
  \cref{sec:taylor_expansions} \\[3pt]
5.4 & semi-important & ``Use the definition of directional derivative,
  $\partial_v f(x) = \left.\frac{\partial}{\partial t}\right|_{t=0} f(x+tv)$. Solve a few of
  these, but not all.'' & \cref{sec:directional_derivatives} \\[3pt]
5.5 & optional & --- & \cref{sec:taylor_expansions} \\[3pt]
5.6 & optional & ``Interesting, but irrelevant for the course'' & \cref{sec:taylor_expansions} \\[3pt]
5.7 & optional & ``Interesting for physicists, will be covered in classical mechanics.'' &
  \cref{sec:chain_rule_along_path} \\
\end{tabular}
\end{center}

% Problems 5.4--5.7 likewise lived in content/exercise-sheets/week-05.tex, which nothing ever
% \input. All four are now placed and solved.

\section{Sheet 6}
\label{sec:sheet_index_6}

% Source: Corsin Nick/Class Notes/Week 6.pdf, p. 1
\begin{center}
\begin{tabular}{@{}l l p{0.38\linewidth} l@{}}
\textbf{Problem} & \textbf{Priority} & \textbf{Note} & \textbf{Now in} \\
\hline
6.1 & unmarked & --- & not transcribed \\[3pt]
6.2 & important & ``This is what we covered on Friday. Find a Lemma in your linear algebra
  notes which relates the determinant and trace to the eigenvalues.'' & \cref{sec:hessian_test} \\[3pt]
6.3 & semi-important & ``By using $\alpha+\beta+\gamma=\pi$ for the angles, you can express
  this as an optimization problem in two variables.'' & not transcribed \\[3pt]
6.4 & unmarked & --- & not transcribed \\[3pt]
6.5 & optional & ``See `data analysis'.'' & not transcribed \\[3pt]
6.6 & optional & ``Do two or three, not all of them.'' & not transcribed \\[3pt]
6.7 & important & --- & \cref{sec:convexity} \\[3pt]
6.8 & important & --- & \cref{sec:hessian_test} \\[3pt]
6.9 & important & --- & \cref{sec:hessian_test} \\[3pt]
6.10--6.12 & important & ``These are good exercises, but how many optimization problems do
  you really need?'' & \cref{sec:lagrange_multipliers}, \cref{sec:hessian_test} \\
\end{tabular}
\end{center}

\begin{ainote}
Sheet 6 is incomplete. Only the problems tagged \textbf{important} (6.2 and 6.7--6.12)
were ever transcribed, and 6.1 and 6.3--6.6 have not been. Everything that \emph{is} transcribed
is solved. This is a gap in the source material, not a casualty of the restructure.

\textbf{6.1 is the exception, and it is deliberate.} It asks for two things about a convex
$f \in C^2(\mathbb{R}^n)$: that a critical point is exactly a global minimiser, and for a
non-negative example whose infimum is not attained. Both are already in the document:
\cref{ex:convexity_characterizations}\textbf{(c)} and
\cref{ex:convex_function_unattained_infimum}, the latter with precisely the required
$f(x) = e^{x_1+x_2}$. Transcribing 6.1 would print the same two questions a second time, so the
row above stays \qt{not transcribed} while the mathematics is covered.
% Considered and rejected 2026-08-11.
\end{ainote}

\section{Sheet 7}
\label{sec:sheet_index_7}

% Source: Corsin Nick/Class Notes/Week 7.pdf, p. 1
\begin{center}
\begin{tabular}{@{}l l p{0.42\linewidth} l@{}}
\textbf{Problem} & \textbf{Priority} & \textbf{Note} & \textbf{Now in} \\
\hline
7.1 & important & ``First think of a diffeomorphism $(a,b) \leftrightarrow \mathbb{R}$ and
  adjust it such that $(a,b)=(0,1)$.'' & \cref{sec:inverse_function_theorem} \\[3pt]
7.2--7.4 & unmarked & --- & not transcribed \\[3pt]
7.5 & important & ``This is a great exercise to help with understanding of IFT.'' &
  \cref{sec:implicit_function_theorem} \\[3pt]
7.6 & important & ``You will need these coordinates very often over the next weeks.'' &
  \cref{sec:tangent_normal_spaces} \\[3pt]
7.7 & semi-important & ``In \textbf{1)}, only the M\"obius strip should be enough.'' &
  not transcribed \\[3pt]
7.8 & semi-important & ``Do \textbf{1)} and \textbf{2)}. The others will be in the lecture.'' &
  not transcribed \\[3pt]
7.9 & unmarked & --- & not transcribed \\
\end{tabular}
\end{center}

% Moved out of the typeset text (the book names tutors only in source citations):
% Corsin's handwritten notes read "In 7)," and "Do 7) and 2)", but neither exercise 7.7 nor 7.8 has
% a sub-item numbered 7). Both are numbered 1.--3. (7.7) and 1.--4. (7.8). Read as 1) above (a
% plausible misreading of his handwritten "1", which this week is slightly closed at the top); this
% also makes the notes self-consistent, since the referenced sub-items (the Moebius strip inside
% 7.7's part 1, and the parametrization/tangent-space parts of 7.8) exist under number 1, not 7.

\begin{ainote}
Like sheet 6, sheet 7 is only partly transcribed. The three tagged \textbf{important},
namely 7.1, 7.5 and 7.6, are in the document and solved. Problems 7.2--7.4 and 7.7--7.9 were
never transcribed.
\end{ainote}

\section{Sheet 8}
\label{sec:sheet_index_8}

% Source: Corsin Nick/Class Notes/Week 8.pdf, p. 1
\begin{center}
\begin{tabular}{@{}l l l@{}}
\textbf{Problem} & \textbf{Priority} & \textbf{Now in} \\
\hline
8.1 & --- & \cref{sec:multivariate_integration} \\[3pt]
8.2 & important & \cref{sec:multivariate_integration} \\[3pt]
8.3 & --- & not transcribed \\[3pt]
8.4 & important & \cref{sec:multivariate_integration} \\[3pt]
8.5 & important & \cref{sec:change_of_variables} \\[3pt]
8.6--8.8 & --- & not transcribed \\
\end{tabular}
\end{center}

% 8.1 was added on 2026-08-11 as ex:jordan_null_products; the row above used to read "not
% transcribed". Sheet 8's other transcribed problems had no solutions at all until the
% restructure (the old Solutions section held only a stale placeholder); all are now solved.
% Standing decision: no more of Sascha Brack's annotated sheet copies are mined. They are copies
% of the official sheet rather than independent notes, so the yield is cross-tutor priority
% agreement plus short margin hints, not worth the reading cost next to Corsin's priority page.
% An ainote here used to say that sheet 8 is the one sheet whose hints partly come from Sascha
% Brack's annotated copy; the hints themselves carry that attribution in their titles, so it was
% removed from the typeset text (the book names tutors only in source citations).

\section{Sheet 9}
\label{sec:sheet_index_9}

% Source: Corsin Nick/Class Notes/Week 9.pdf, p. 1
Only \textbf{9.1} and \textbf{9.6} carry an explicit colour marker on the priority page;
\textbf{9.2} is annotated with a hint but no clearly distinguishable colour, so it is listed as
unmarked below even though a hint is quoted for it. \textbf{9.3--9.5} and \textbf{9.7} carry
neither colour nor hint.

\begin{center}
\begin{tabular}{@{}l l p{0.40\linewidth} l@{}}
\textbf{Problem} & \textbf{Priority} & \textbf{Note} & \textbf{Now in} \\
\hline
9.1 & important & (hint on coordinates, quoted with the problem) &
  \cref{sec:volume_embedded_surfaces} \\[3pt]
9.2 & unmarked & (hint on the domain of part 2, quoted with the problem) &
  \cref{sec:length_of_curve} \\[3pt]
9.3--9.5 & unmarked & --- & not transcribed \\[3pt]
9.6 & semi-important & ``This may become important in later courses.
  $\langle\cdot,\cdot\rangle$ is the standard inner prod.'' & not transcribed \\[3pt]
9.7 & unmarked & --- & not transcribed \\
\end{tabular}
\end{center}

\begin{ainote}
Only 9.1 is tagged important, but its companion 9.2 carries an explicit worked hint (unlike 9.3--9.5 and 9.7, which carry nothing at all), so both are reproduced and solved.
9.3--9.7 were never transcribed.
\end{ainote}

\section{Sheet 10}
\label{sec:sheet_index_10}

% Source: Corsin Nick/Class Notes/Week 10.pdf, p. 1
\begin{center}
\begin{tabular}{@{}l l p{0.38\linewidth} l@{}}
\textbf{Problem} & \textbf{Priority} & \textbf{Note} & \textbf{Now in} \\
\hline
10.1 & important & ``Use spherical coordinates with $r\equiv1$.'' &
  \cref{sec:volume_embedded_surfaces} \\[3pt]
10.2 & important & (worked parametrization of the torus, quoted with the problem) &
  \cref{sec:volume_embedded_surfaces} \\[3pt]
10.3 & important & ``1.)--4.), 5.)'' & \cref{sec:length_of_curve} \\[3pt]
10.4--10.6 & semi-important & --- & not transcribed \\
\end{tabular}
\end{center}

\begin{ainote}
For \textbf{10.3}, the note \qt{1.)--4.), 5.)} appears to flag that the starred sub-part
\textbf{5.} (the harder \qt{if and only if} characterization of a bounded $C^1$ domain, marked
(*) on the official sheet) is worth attempting alongside the more routine parts 1--4, not
skipped as the (*) marker might suggest. 10.4--10.6 (tractrix, logarithmic spiral, moment of
inertia of an ellipse) were never transcribed.
\end{ainote}

\section{Sheet 11}
\label{sec:sheet_index_11}

% Source: Corsin Nick/Class Notes/Week 11.pdf, p. 1
\begin{center}
\begin{tabular}{@{}l l l@{}}
\textbf{Problem} & \textbf{Priority} & \textbf{Now in} \\
\hline
11.1 & important & \cref{sec:vector_calculus_operators} \\[3pt]
11.2--11.4 & semi-important & not transcribed \\[3pt]
11.5 & important & \cref{sec:flux_divergence_theorem} \\[3pt]
11.6--11.9 & optional & not transcribed \\
\end{tabular}
\end{center}

\begin{ainote}
Problems 11.2--11.4 and 11.6--11.9 were never transcribed, consistent with the standing
decision to follow the priority marking. Between them they cover improper-integral
classification, a standard flux computation, a cylinder's fluxes and volumes, the heat kernel,
a divergence-theorem reality check, a multiple-choice question on Lipschitz continuity, and
differential-form computations.
\end{ainote}

\section{Sheet 12}
\label{sec:sheet_index_12}

% Source: Corsin Nick/Class Notes/Week 12.pdf, p. 1
\begin{center}
\begin{tabular}{@{}l l l@{}}
\textbf{Problem} & \textbf{Priority} & \textbf{Now in} \\
\hline
12.1 & semi-important & \cref{sec:greens_theorem_conservative_fields} \\[3pt]
12.2 & optional & not transcribed \\[3pt]
12.3 & important & \cref{sec:area_via_green} \\[3pt]
12.4 & important & \cref{sec:area_via_green} \\[3pt]
12.5 & important & \cref{sec:area_via_green} \\[3pt]
12.6--12.7 & semi-important & not transcribed \\[3pt]
12.8 & important & \cref{sec:area_via_green} \\[3pt]
12.9 & important & \cref{sec:integration_of_forms} \\
\end{tabular}
\end{center}

\begin{ainote}
12.2 (existence of a potential via a path integral, building the tool used implicitly in
\cref{ex:conservative_vector_field_potential}), 12.6 (a Stokes reality check structurally identical to \cref{ex:stokes_work_along_paths} part
\textbf{(c)}) and 12.7 (four standard $\curl$/$\divg$ product-rule identities) have not been
transcribed. 12.1 is \cref{ex:path_concatenation_work}, placed beside \cref{def:line_integral}
rather than with the other sheet-12 problems, because it is a property of that definition and is
used implicitly from there onwards.
\end{ainote}

\section{Sheet 13}
\label{sec:sheet_index_13}

% Correction: Claude Opus 5 (High), 2026-08-11. This section previously asserted that "whether a
% thirteenth sheet was issued is not recorded in any source available here". That was false:
% exercises/Ex13_Analysis2_eng.pdf sits in the same directory as the other twelve, and the FS 2024
% edition ships its own sheet 13 as well (exercises_2024/ps13.pdf). The claim was presumably
% inherited from the week-13 chapter, which carried no \exercisesheet heading, and never rechecked
% against the directory --- an instance of the "a grep is not a survey" trap in gemini.md. Found
% by the user.

% Source: exercises/Ex13_Analysis2_eng.pdf
\begin{center}
\begin{tabular}{@{}l p{0.46\linewidth} l@{}}
\textbf{Problem} & \textbf{What it asks} & \textbf{Now in} \\
\hline
13.1 & An inhomogeneous second-order initial value problem, solved and its solution shown
  unique & not transcribed \\[3pt]
13.2 & Three constant-coefficient homogeneous equations, and what changes over $\mathbb{C}$
  & not transcribed \\[3pt]
13.3 & Three non-homogeneous problems, the last by the integrating-factor trick
  & not transcribed \\[3pt]
13.4 & Four equations put in normal form, tested against Cauchy--Lipschitz, then solved
  & \cref{sec:cauchy_lipschitz} \\[3pt]
13.5 & Multiple choice: when the solution set is a vector space; bounded solutions of
  $\dot x = Ax$; convergence of the Picard iteration & not transcribed \\[3pt]
13.6 & True or false & not transcribed \\
\end{tabular}
\end{center}

% The week-13 chapter carried no exercise-sheet heading, and until 2026-08-11 this appendix
% recorded that no thirteenth sheet was known to exist; see the correction comment above.
\begin{ainote}
Sheet 13 is the course's ODE sheet, and its material belongs to \cref{ch:odes} rather than to any
chapter of the main text.

Only 13.4 is transcribed, as \cref{ex:serra_normal_form_cauchy_lipschitz}. It was chosen because it
is the one problem on the sheet that is not a computation: it asks in each case \emph{whether} the
existence and uniqueness theorem applies before asking for a solution, which is the question
\cref{sec:cauchy_lipschitz} exists to answer. 13.1--13.3 are solution techniques that
\cref{sec:solving_linear_systems_odes} and \cref{sec:inhomogeneous_odes} already cover, and 13.5
asks about the long-time behaviour of $\dot x = Ax$, which this document does not develop, and
about convergence of the Picard iteration, which \cref{sec:picard_iteration} treats only briefly.
\end{ainote}
% Correction: Claude Opus 5.5 (Extra) --- 13.5 was said to reach into "the Picard iteration ...
% which this document does not develop", but sec:picard_iteration exists; and the closing
% sentence called ex:area_boundary_functional and ex:area_ellipse "exercises", which they stopped
% being when they were reclassified as worked examples (see 26-stokes/05). Removed: it answered
% nothing about sheet 13.
